\SourceLeaf{347}%
\par\indent When an investigation involves many directions of straight lines in space, it helps to represent the directions by points on the surface of a fixed sphere. These points are the endpoints of radii drawn parallel to the respective lines. The center and radius of this \textit{auxiliary sphere} are entirely arbitrary. The unit of length can serve as its radius. This is essentially the method used constantly in astronomy, where all directions are referred to an imaginary celestial sphere of infinitely large radius. Spherical trigonometry and several other theorems can then solve the problems that arise in comparing the various directions. The author has added a new theorem to these that can often be applied.

If one represents the direction of the normal erected at each point of a curved surface by the corresponding point on the sphere, using the method just described, then each point on the curved surface has a corresponding point on the surface of the auxiliary sphere. Generally speaking, a line on the curved surface then corresponds to a line on the auxiliary sphere, and a patch of the former surface to a patch of the latter. The less a patch departs from a plane, the smaller its corresponding patch on the sphere. The area of the corresponding spherical patch therefore suggests itself as a measure of the total curvature to be assigned to the patch of the curved surface. The author accordingly calls this area the \textit{total curvature} of the corresponding patch of the curved surface. Along with size, the \textit{orientation} of the parts also matters. Regardless of their relative sizes, the two patches may have similar or reversed orientations. A positive or negative sign before the total curvature can distinguish these cases. The distinction has a definite meaning only if the figures are taken on specified sides of the two surfaces. The author takes the outside of the sphere and the side of the curved surface on which the normal is erected. It follows that the sign is positive for convex-convex or concave-concave surfaces, which are essentially alike, and %
